\documentclass[a4paper,11pt]{article}
\usepackage[utf8]{inputenc}
\usepackage[a4paper]{geometry}
\geometry{top=2cm, bottom=2cm, left=2.3cm, right=1.9cm}
\usepackage[T1]{fontenc}
\usepackage[english]{babel}
\usepackage{amsmath}
\usepackage{amssymb}
\usepackage{mathtools}
\usepackage{enumerate}
\usepackage{amsthm}
\usepackage{multirow}
\usepackage{gensymb}
\usepackage{diagbox}
\usepackage{enumitem}
\usepackage{graphicx}
\usepackage{color}
\usepackage{xfrac}
\usepackage{setspace}
\usepackage{hyperref}
\usepackage{array}
\usepackage{physics}
\usepackage[nameinlink]{cleveref}
\usepackage{titlesec}
\usepackage{authblk}
\usepackage[most]{tcolorbox}
\usepackage{listings}
\usepackage{xcolor}

\onehalfspacing
\linespread{1.0}
\definecolor{amethyst}{rgb}{0.6, 0.4, 0.8}
\definecolor{green}{rgb}{0.0, 1.0, 0.0}
\definecolor{leanblue}{RGB}{0,0,160}
\definecolor{leangreen}{RGB}{0,120,0}
\definecolor{leanred}{RGB}{160,0,0}
\definecolor{leancomment}{RGB}{120,120,120}
\hypersetup{colorlinks, citecolor=green, linkcolor=blue, urlcolor=blue}
\setlength{\parindent}{16pt}

\lstdefinelanguage{Lean}{
	keywords={
		theorem, lemma, def, structure, inductive, class,
		namespace, section, variable, variables,
		open, import, universe, universes,
		begin, end, match, with, if, then, else,
		by, fun, assume, have, show, from,
		forall, exists, let, in
	},
	sensitive=true,
	comment=[l]{--},
	morecomment=[s]{/-}{-/},
	string=[b]",
}

\lstset{
	language=Lean,
	basicstyle=\ttfamily\small,
	keywordstyle=\color{leanblue}\bfseries,
	commentstyle=\color{leancomment}\itshape,
	stringstyle=\color{leanred},
	numbers=left,
	numberstyle=\tiny,
	mathescape=true,
	stepnumber=1,
	frame=single,
	breaklines=true,
	tabsize=2,
	showstringspaces=false,
	captionpos=b
}

\lstset{
	literate=
	{ℤ}{{$\mathbb{Z}$}}1
	{∀}{{$\forall$}}1
	{∃}{{$\exists$}}1
	{→}{{$\to$}}1
	{∧}{{$\land$}}1
	{∈}{{$\in$}}1
	{≤}{{$\leqslant$}}1
	{≥}{{$\geqslant$}}1
	{⟨}{{$\langle$}}1
	{⟩}{{$\rangle$}}1
	{⁻¹}{{$^{-1}$}}1
	{⟨}{{$\langle$}}1
	{⟩}{{$\rangle$}}1
	{↔}{{$\Leftrightarrow$}}1
	{∧}{{$\wedge$}}1
	{ℝ}{{$\R$}}1
	{⇒}{{$\Rightarrow$}}1
	{δ}{{$\delta$}}1
	{K₂}{{$K_2$}}1
	{K₁}{{$K_1$}}1
	{x₀}{{$x_0$}}1
	{ε}{{$\varepsilon$}}1
	{≠}{{$\neq$}}1
	{·}{{\textperiodcentered}}1
}

\crefname{lstlisting}{Listing}{Listings}
\Crefname{lstlisting}{Listing}{Listings}

\theoremstyle{plain}
\newtheorem{theorem}{Theorem}
%\newtheorem{lemma}[theorem]{Lemma}
\newtheorem{definition}[theorem]{Definition}

\newcommand\restr[2]{{% we make the whole thing an ordinary symbol
		\left.\kern-\nulldelimiterspace % automatically resize the bar with \right
		#1 % the function
		\littletaller % pretend it's a little taller at normal size
		\right|_{#2} % this is the delimiter
}}
\newcommand{\R}{\ensuremath{\mathbb{R}}} % set of real numbers
\newcommand{\littletaller}{\mathchoice{\vphantom{\big|}}{}{}{}}

\tcbuselibrary{theorems}

% tcolorbox lemma environment with label prefix 'lem'
\newtcbtheorem[number within=section]{lemma}{Lemma}{
	fonttitle=\bfseries,
	coltitle=black,
	colback=gray!5,
	colframe=gray!60,
	boxrule=0.5pt,
	arc=2pt,
	left=6pt,
	right=6pt,
	top=6pt,
	bottom=6pt
}{lem}

% Tell cleveref how to print references to lemmas
\crefname{lemma}{lemma}{lemmas}
\Crefname{lemma}{Lemma}{Lemmas}

% --------------------
% Title formatting
% --------------------
\titleformat{\section}
{\normalfont\Large\bfseries}{\thesection}{1em}{}


\begin{document}
	\sffamily
	%\title{\textbf{Formalising Mathematics} \\ \textbf{Project 1} \\ \textbf{Finite Jensen inequality for convex functions on $\R$}}
	%\author{Joan Arenillas i Cases, \ 06067859}
	%\maketitle
	
	% --------------------
	% Title
	% --------------------
	\begin{center}
		{\LARGE\bfseries
			Formalising Mathematics: \\[0.3em]
			Proving basic notions of limits and boundedness\\[0.3em]
			\par}
	\end{center}
	
	\vspace{0.4cm}
	
	% --------------------
	% Author
	% --------------------
	\begin{center}
		\textbf{Joan Arenillas i Cases}\textsuperscript{1}
	\end{center}
	
	\vspace{0.5cm}
	
	\noindent
	\textsuperscript{1}MSc in Applied Mathematics, Imperial College London.\\
	CID: 06067859, \ \href{mailto:ja1225@ic.ac.uk}{ja1225@ic.ac.uk}
	
	\vspace{0.6cm}
	
	% --------------------
	% Keywords
	% --------------------
	\vspace{0.5cm}
	
	% --------------------
	% Abstract
	% --------------------
	\noindent
	\textbf{Abstract.}
	In this report I state elemental definitions of limits of functions and prove that if a function $f$ has a limit at some point $x_0$ of its domain, then $f$ is necessarily bounded on some neighbourhood of $x_0$. In the way to prove this first-year result, we define concepts on limits, working through some examples. In the second part of the report we formalise our main result, as well as other auxiliary---but necessary---definitions. At the end, we discuss the challenges faced, reflecting on the assumptions.
	
	\vspace{0.5cm}
	
	\section{Introduction}
	
	The goal of this report is to formalise in Lean a basic result in real analysis. It states the following:\\
	
	\begin{lemma*}{Boundedness around a point}{}
			Let $f:(a,b)\rightarrow \R$ for some $a,b\in\R$. If there exists the limit of $f$ at $x_0\in(a,b)$, then $f$ is bounded on some neighbourhood of $x_0$.
	\end{lemma*}

	
	We will see in \Cref{sec:informal} that this result can be proved in a few lines with a small number of definitions on limits and boundedness. However, we will discover in \Cref{sec:formal} that the Lean implementation is far more intricate. It will require handling the proof by cases and considering intervals and bounds very carefully. We will ponder over the challenges we faced in \Cref{sec:challenges}, and we will propose some ways to move forward in \Cref{sec:conclu}.
	
	There are countless sources explaining the basic properties of real analysis in one variable. Unless otherwise stated, in this report I have followed \cite{manyosas}.

	\section{Informal Proof}\label{sec:informal}
	
	To reach the proof of our lemma, we must first define some concepts. Let $f: (a, b) \rightarrow \R$ be a function and consider $x_0\in(a, b)$. We say that $\ell$ is the \emph{limit of the function $f$ at $x_0$} and we write $\lim_{x\rightarrow x_0}f(x)=\ell$ if for all $\varepsilon>0$ there exists $\delta>0$ such that $\abs{f(x)-\ell}<\varepsilon$ if $0 < \abs{x-x_0} < \delta$. Observe that here we are always supposing that the limit $\ell$ is finite and that $f$ does not take infinite values in $(a,b)$.
	
	Although we will not need this result to prove the desired Lemma we will also introduce the notion of continuity to practise with Lean definitions. Specifically, we will say that the function $f$ \emph{is continuous at $x_0\in(a,b)$} if the limit $\lim_{x\rightarrow x_0}f(x)$ exists and is equal to $f(x_0)$. Naturally, we will say that $f$ \emph{is continuous on $(a,b)$} if $f$ is continuous at $x_0$ for every $x_0\in(a,b)$. Lastly, we say that $f:\R\rightarrow\R$ is \emph{continuous on $\R$} if it is continuous at every $x_0\in\R$.
	
	As an example, consider $g(x)=x$ defined for all $x\in\R$. It is now easy to check that $g(x_0)=\lim_{x\rightarrow x_0} g(x)=x_0$ for any $x_0 \in \R$. Indeed, fix some $\varepsilon>0$ and define $\delta:=\varepsilon>0$. Now, for any $x\in\R$ with $0 < \abs{x-x_0} < \delta=\varepsilon$ we have that $\abs{g(x)-x_0}=\abs{x-x_0}<\varepsilon$.
	
	It is also worth noting that it makes sense to talk about ``the'' limit of a function $f$ because if it exists, this limit is unique (we will not formally nor informally prove this result).
	
	We will now introduce the fundamental concept for the Lemma we aim to prove in Lean. Let $I$ be an interval and consider $f : I \rightarrow \R$. We say that $f$ is \emph{bounded on $I$} if there exist some constants $m,M\in\R$ such that $m \leqslant f(x) \leqslant M$ for all $x \in I$. This condition can be equivalently stated as follows.\\
	
	\begin{lemma}{Equivalence of boundedness definitions}{}
		A function $f : I \rightarrow \R$ is bounded on $I$ if and only if there exists some real constant $K\geqslant0$ such that $\abs{f(x)} \leqslant K$ for every $x \in I$. 
	\end{lemma}
	\begin{proof}
		If $f$ is bounded on $I$ then there exist some real constants $m$ and $M$ such that $m \leqslant f(x) \leqslant M$ for all $x \in I$, so that $\abs{f(x)}\leqslant\max\{\abs{m},\abs{M}\}=:K$. Conversely, if $\abs{f(x)}\leqslant K$, $\forall x \in I$ for some $K\geqslant0$, we deduce that $-K \leqslant f(x) \leqslant K$ for all $x \in I$.
	\end{proof}
	
	In a similar way, we say that $f:I\rightarrow\R$ is bounded on $J\subseteq I$ if its restriction $\restr{f}{J}$ is bounded on $J$.
	
	Observe that the existence of the limit of a function in one point $x_0\in I$ does not provide us with global information about the function, but rather gives information on a neighbourhood of $x_0$. We now state and prove the main result we are interested in, which gives information on the function $f$ near $x_0$.\\
	
	\begin{lemma}{Boundedness around a point}{boundedness}
		\label{lem:boundedness}
		Let $f:(a,b)\rightarrow \R$ for some $a,b\in\R$. If there exists the limit of $f$ at $x_0\in(a,b)$, then $f$ is bounded on some neighbourhood of $x_0$.
	\end{lemma}
	\begin{proof}
		Let $\ell:=\lim_{x\rightarrow x_0}f(x)$. Let us apply the definition of limit with $\varepsilon=1$. Then, we have that there exists some $\delta>0$ such that $\abs{f(x)- \ell} < 1$ if $0 < \abs{x-x_0} < \delta$. Then, by the triangle inequality we have that $\abs{f(x)}= \abs{(f(x)-\ell) + \ell}\leqslant \abs{f(x)-\ell}+\abs{\ell}< 1 + \abs{\ell}$ if $0 < \abs{x-x_0} < \delta$. To conclude, let $K:=\max\{1 + \abs{\ell}, \abs{f(x_0)}\}<\infty$ and note that $K\geqslant0$. We then deduce that $\abs{f(x)}\leqslant K$ if $x \in (x_0-\delta, x_0 + \delta)$, so that $f$ is bounded on this neighbourhood of $x_0$.
	\end{proof}
	
	\section{Formal Proof}\label{sec:formal}
	 
	 We now turn our attention to formalising \Cref{sec:informal} in Lean\footnote{I acknowledge that I have used AI tools to check the spelling and grammatical structure of this report. Also, GitHub Copilot has been used to suggest possible tactics and lemmas in the proofs.}. We are going to follow the structure of the previous section. 
	 
	 \subsection{Preliminary definitions and examples}
	 
	 We start by defining the interval endpoints with \lstinline|variable (a b : $\R$)| and further defining the notion of limit of a function $f:(a,b)\rightarrow \R$ at a point $x_0$ (with limit $\ell$). \\
	
	\begin{lstlisting}[caption={Definition of limit of $f$ at $x_0$.}]
		def limit_f_x_0 {a b : ℝ} (l x₀ : ℝ) (f : Set.Ioo a b → ℝ) : Prop :=
			∀ ε > 0, ∃ δ > 0, ∀ x : Set.Ioo a b, |x - x₀| < δ ∧ |x - x₀| > 0 → |f x - l| < ε
	\end{lstlisting}
	
	Observe that the open interval is defined through \lstinline|Set.Ioo a b|, where we do not need to specify the condition $a<b$ because Lean will consider the empty set if $a>b$. We continue by giving the definition of continuity of $f$ at $x_0$, continuity on the whole interval $(a,b)$ and continuity on the real line.\\
	
	\begin{lstlisting}[caption={Definition of continuity at $x_0$.}, label={list:1}]
		def continuous_x0 {a b : ℝ} (f : Set.Ioo a b → ℝ) (x₀ : ℝ) (hx₀ : x₀ ∈ Set.Ioo a b) : Prop :=
			limit_f_x_0 (l := f ⟨x₀, hx₀⟩) x₀ f
	\end{lstlisting}
	
	\begin{lstlisting}[caption={Definition of continuity on $(a,b)$.}, label={list:2}]
		def continuous_on_interval (f : Set.Ioo a b → ℝ) : Prop :=
			∀ x₀ (hx₀ : x₀ ∈ Set.Ioo a b), continuous_x0 f x₀ hx₀
	\end{lstlisting}
	
	\begin{lstlisting}[caption={Definition of continuity on all $\R$.}, label={list:3}]
		def continuous_on_R (f : ℝ → ℝ) : Prop :=
			∀ a b x₀ (hx₀ : x₀ ∈ Set.Ioo a b),
				continuous_x0 (fun x : Set.Ioo a b => f x) x₀ hx₀
	\end{lstlisting}
	
	The definitions in \crefrange{list:1}{list:3} are very standard and closely follow the definitions we introduced in \Cref{sec:informal}. It is worth noting, however, the use of \lstinline|{a b : $\R$}| at the beginning of the definitions of limit and continuity at $x_0$. By doing this, we will not need to specify $a$ and $b$ as arguments when using these definitions later in the code.
	
	It is also worth noting that we defined continuity over $\R$ (\cref{list:3}) by using the notion of continuity at a point $x_0$ we had previously introduced. This forces us to define continuity on $\R$ by saying that a continuous function on $\R$ is a continuous function at $x_0$ for all $x_0\in(a,b)$ and for all $a,b\in\R$. To formalise this definition, we introduced a new variable $x_0$ and the restriction of $f:\R\rightarrow\R$ to the interval $(a,b)$. In particular, the function \lstinline|f| passed to \lstinline|continuous_x0| is restricted to $(a,b)$ using the equivalent of a $\lambda$-function definition.
	
	Observe that we need this non-standard definition of continuity on $\R$ to be able to use the notion of continuity defined in \lstinline|continuous_x0|, which expects a function defined in some interval $(a,b)$. Obviously, Lean can handle continuity on $\R$ more cleanly by using the Mathlib definition \lstinline|Continuous f| \cite{LeanContinuity}.
	
	Next, we show that the function $g:\R\rightarrow\R$ given by $x\mapsto x$ has limit equal to $x_0$ at every point $x_0\in\R$. We define this function via \lstinline|def linear_fun : ℝ → ℝ := fun x => x|.\\
	
	\begin{lstlisting}[caption={Example that $\lim_{x\rightarrow x_0}x=x_0$.}, label={list:4}]
		example : continuous_on_R linear_fun := by
			-- introduce assumptions that `continuous_on_R` expects.
			intro a b x₀ hx₀
			-- unfold continuity (and limit) at x₀
			dsimp [continuous_x0, limit_f_x_0]
			intro ε hε
			-- choose δ = ε and provide proof that δ > 0
			refine ⟨ε, hε, ?_⟩
			-- hx is the hypothesis: |x - x₀| < ε and |x - x₀| > 0
			intro x hx
			-- `[linear_fun]` replaces `f` by its specific definition
			simpa [linear_fun] using hx.1
	\end{lstlisting}
	
	There are two tactics worth explaining in \cref{list:4}. The first one is \lstinline|dsimp|, which unfolds the definition of continuity on $\R$. This means writing out the definition of continuity at $x_0$ and, next, the definition of limit at $x_0$. The second tactic is \lstinline|	refine ⟨ε, hε, ?_⟩|, which changes the goal by using $\delta:=\varepsilon$ in the definition of limit at $x_0$. 
	
	\subsection{Definitions of boundedness}
	
	We now state the two definitions of boundedness on an interval in \cref{list:5} and prove their equivalence in \Cref{list:6}.\\
	
	\begin{lstlisting}[caption={Two equivalent definitions of boundedness.}, label={list:5}]
		def bounded_on_interval_1 (f : Set.Ioo a b → ℝ) : Prop :=
			∃ (M m : ℝ), ∀ x : Set.Ioo a b, m ≤ f x ∧ f x ≤ M
		
		def bounded_on_interval_2 (f : Set.Ioo a b → ℝ) : Prop :=
			∃ K : ℝ, 0 ≤ K ∧ ∀ x : Set.Ioo a b, |f x| ≤ K
	\end{lstlisting}
	
	\begin{lstlisting}[caption={The two definitions of boundedness are equivalent.}, label={list:6}]
		lemma bounded_equiv (f : Set.Ioo a b → ℝ) :
			bounded_on_interval_1 a b f ↔
			bounded_on_interval_2 a b f := by
			-- separate the two implications
			constructor
			· -- (Def 1) → (Def 2)
				intro h
				-- unpack hypothesis and substitute
				rcases h with ⟨M, m, hMm⟩
				-- define bound for Def 2 with m, M
				let K0 : ℝ := max (|m|) (|M|)
				-- unpack what needs to be proved with K0
				refine ⟨K0, ?_, ?_⟩
					· -- prove K0 ≥ 0
					have hm : 0 ≤ |m| := abs_nonneg m
					have hmk : |m| ≤ K0 := le_max_left _ _ -- the underscores tell Lean to infer the arguments
					exact hm.trans hmk
					· -- prove |f x| ≤ K0
					intro x
					-- prove -K0 ≤ f x
					have hx_low : -K0 ≤ f x := by
						-- use -K0 ≤ -|m| ≤ m ≤ f x
						have h1 : -K0 ≤ -|m| := by
							have : |m| ≤ K0 := le_max_left _ _
							exact neg_le_neg this -- `this` refers to the last `have` statement
						have h2 : -|m| ≤ m := by linarith [neg_le_abs m]
						have h3 : m ≤ f x := by simpa using (hMm x).1
						exact h1.trans (h2.trans h3)
					-- prove fx ≤ K0
					have hx_up : f x ≤ K0 := by
						-- f x ≤ M ≤ |M| ≤ K0
						have hfxM : f x ≤ M := (hMm x).2
						have hMabs : M ≤ |M| := le_abs_self M
						have habsK : |M| ≤ K0 := le_max_right _ _
						exact hfxM.trans (hMabs.trans habsK)
					exact (abs_le.mpr ⟨hx_low, hx_up⟩) -- mpr uses the reverse implication of `abs_le`
					-- the ⟨,⟩ brackets are used to construct a hypothesis of type A ∧ B
			
			· -- (Def 2) → (Def 1)
					intro h
					rcases h with ⟨K, Kpos, hK⟩
					-- take m = -K, M = K, so, from |f x| ≤ K → -K ≤ f x ≤ K
					refine ⟨K, -K, ?_⟩
					intro x
					-- abs_le.mp: |f x| ≤ K means -K ≤ f x ∧ f x ≤ K
					simpa using (abs_le.mp (hK x)) -- mp uses the direct implication of `abs_le`
	\end{lstlisting}
	
	There are a few tactics worth mentioning in \cref{list:6}:
	\begin{enumerate}
		\item In line 13, the \lstinline|refine| tactic is used in a different way. Here, this tactic creates two new goals associated with \lstinline|K0|: prove that \lstinline|K0| $\geqslant 0$ and $\abs{f(x)}\leqslant$ \lstinline|K0|. For this latter task, we break down the inequality in small pieces, which we prove in dedicated \lstinline|have| statements.
		\item In line 16, the underscores in \lstinline|le_max_left _ _| tell Lean to infer the arguments for the lemma \lstinline|le_max_left|, which states that $a\leqslant \max\{a,b\}$, for $a,b$ in any linearly ordered type.
		\item Throughout this snippet of code we extensively use the \lstinline|trans| lemma, which gives a proof of the transitive property of non-strict inequalities.
		\item In line 36 note how \lstinline|.mpr| obtains the reverse implication in an ``if and only if statement'', while \lstinline|.mp| obtains the direct implication (line 46).
		\item Finally, we use numerous lemmas regarding basic properties of the absolute value, which have been found in LeanSearch \cite{LeanSearch}, using the \lstinline|exact?| and \lstinline|try?| tactics, or with GitHub copilot.
	\end{enumerate}
	
	Next, we introduce the notion of boundedness in a subset $J\subseteq I$ as follows:\\
	
	\begin{lstlisting}[caption={Boundedness of $f:J\subseteq I\to\R$.}, label={list:7}]
		def bounded_on_J {a b : ℝ} (J : Set ℝ) (f : Set.Ioo a b → ℝ) : Prop :=
			∃ K : ℝ, 0 ≤ K ∧ ∀ x ∈ J, ∀ hx : x ∈ Set.Ioo a b, |f ⟨x, hx⟩| ≤ K
	\end{lstlisting}
	
	\subsection{Main Lemma}

	We can now prove the main Lemma of this project:\\
	
	\begin{lstlisting}[caption={A function with limit at $x_0$ is bounded on some interval around $x_0$.}, label={list:8}]
		lemma bounded_if_limit_exists (f : Set.Ioo a b → ℝ) (x₀ l : ℝ) (hx0ab : x₀ ∈ Set.Ioo a b) :
				limit_f_x_0 l x₀ f → ∃ δ > 0, bounded_on_J (Set.Ioo (x₀ - δ) (x₀ + δ)) f := by
			intro hlim
			-- use definition of limit with ε = 1
			obtain ⟨δ, hδpos, hδ⟩ := hlim 1 (by norm_num)
			-- define one possible bound
			let K₁ : ℝ := 1 + |l|
			-- |f x| < K₁ if |x - x₀| < δ and |x - x₀| > 0
			have hδ_to_K1 :
			∀ x : Set.Ioo a b,
				(|x - x₀| < δ ∧ |x - x₀| > 0) → |f x| < K₁ :=
			by
				intro x hx
				-- use triangle inequality to see |f x| ≤ |f x - l| + |l|
				have triang : |f x| ≤ |f x - l| + |l| := by
					-- norm_add_le does: |a + b|  ≤ |a| + |b|
					have := norm_add_le (f x - l) l
					simpa using this
				-- |f x - l| < 1
				have hfxl : |f x - l| < 1 := hδ x hx
				-- transform to: |f x - l| + |l| < 1 + |l| = K₁
				have add_l : |f x - l| + |l| < K₁ := by
					-- add_lt_add_right does: if a < b then a + c < b + c
					simpa [K₁] using add_lt_add_right hfxl |l|
				-- conclude |f x| < K₁
				exact lt_of_le_of_lt triang add_l
			
			-- define the final upper bound
			let K₂ : ℝ := max K₁ |f ⟨x₀, hx0ab⟩|
			-- simplify hypothesis and split goals
			refine ⟨δ, hδpos, ?_⟩
			refine ⟨K₂, ?_, ?_⟩
			· -- prove 0 ≤ K₂
				-- prove 1 ≤ K₁ = 1 + |l|, since |l| ≥ 0
				have h1leqK1 : 1 ≤ K₁ := by
					have habs : 0 ≤ |l| := abs_nonneg l
					-- 1 ≤ 1 + |l|
					linarith [habs]
				-- prove K₁ ≤ K₂ = max K₁ |f x₀|
				have hK1leqK2 : K₁ ≤ K₂ := by
					simpa [K₂] using le_max_left K₁ |f ⟨x₀, hx0ab⟩|
				exact (le_trans (by norm_num) (h1leqK1.trans hK1leqK2))
			· -- prove |f x| ≤ K₂ for x ∈ (x₀ - δ, x₀ + δ)
				intro x hxI hxDom
				by_cases hxeq : x = x₀ -- observe x satisfies both hxI and hxDom
				· -- bound the case x = x₀
					have hx0_le_K2 : |f ⟨x₀, hx0ab⟩| ≤ K₂ :=
						le_max_right _ _
					have : |f ⟨x, hxDom⟩| ≤ K₂ := by
						simpa [hxeq] using hx0_le_K2
					exact this
				· -- bound the case x ≠ x₀, i.e., |x - x₀| < δ and 0 < |x - x₀|
					have hx_abs_lt : |x - x₀| < δ := by
						have hL : -δ < x - x₀ := by linarith [hxI.1]
						have hR : x - x₀ < δ := by linarith [hxI.2]
						exact abs_lt.mpr ⟨hL, hR⟩
					have hpos_abs : |x - x₀| > 0 := abs_pos.mpr (sub_ne_zero.mpr hxeq)
					-- bound |f x| < K₁ for x ∈ 0 < |x - x₀| < δ
					have flK1 : |f ⟨x, hxDom⟩| < K₁ :=
						-- use |x - x₀| < δ and 0 < |x - x₀|
						hδ_to_K1 ⟨x, hxDom⟩ ⟨by simpa using hx_abs_lt, hpos_abs⟩
					-- K₁ ≤ K₂
					have K₁_le_K₂ : K₁ ≤ K₂ := by
						simpa [K₂] using (le_max_left _ _)
					-- conclude |f x| ≤ K₂
					exact (le_of_lt flK1).trans K₁_le_K₂
	\end{lstlisting}
	
	The tactics and lemmas we use in \cref{list:8} are very similar to the ones we described in \cref{list:6}. In the following lines I will outline the general structure of the main Lemma's proof.
	\begin{enumerate}
		\item We use the definition of limit with $\varepsilon:=1$ and use the second definition of boundedness. Using this definition, we first need to prove that $\abs{f(x)}<K_1:=1+\abs{\ell}$ for all $x\in(x_0-\delta, x_0+\delta)$ except for $x=x_0$.
		\item We prove this by first seeing that $\abs{f(x)}\leqslant \abs{f(x)-\ell}+\abs{\ell}$ (through the triangle inequality), next using that $\abs{f(x)-\ell}<1$, and finally building the sequence $\abs{f(x)}\leqslant\abs{f(x)-\ell}+\abs{\ell}<1+\abs{\ell}=K_1$.
		\item Since we also need to consider the point $x_0$, we define $K_2:=\max\{K_1, \abs{f(x_0)}\}$, and prove that $K_2$ is the non-negative constant that satisfies the definition of boundedness in $(x_0-\delta, x_0+\delta)$.
		\item For this goal, we need to prove two goals (through the \lstinline|refine| tactic), the first of which being to prove that $K_2\geqslant 0$. We do this by proving that $1\leqslant K_1=1+\abs{\ell}$ and that $K_1\leqslant K_2=\max\{K_1, \abs{f(x_0)}\}$.
		\item The second and last goal is to prove that $\abs{f(x)}\leqslant K_2$ for all $x\in(x_0-\delta, x_0+\delta)$. The \lstinline|by_cases| tactic allows us to separate the cases $x=x_0$ and $x\neq x_0$. The first case is straightforward, and the second requires showing that $\abs{f(x)}<K_1\leqslant K_2$.
	\end{enumerate}
	
	Finally, we run the \lstinline|#lint| command to check for formatting, and it shows that all the linting tests have passed, and no errors occur.
	
	\section{Challenges faced}\label{sec:challenges}
	
		I found it particularly tricky to understand the correct use of the symbols $``\{\}"$ and $``()"$ when used as hypotheses. This was particularly noticeable in the definition of the limit of a function, where the following definition failed when called inside the main Lemma.\\
		
		\begin{lstlisting}[caption={(Incorrect) definition of limit at $x_0$.}]
			def limit_f_x_0 (a b : ℝ) (l x₀ : ℝ) (f : Set.Ioo a b → ℝ) : Prop :=
				∀ ε > 0, ∃ δ > 0, ∀ x : Set.Ioo a b, |x - x₀| < δ ∧ |x - x₀| > 0 → |f x - l| < ε
		\end{lstlisting}
		
		However, changing \lstinline|(a b : $\R$)| to curly brackets \lstinline|{a b : $\R$}| makes all the difference, since this produces no ``application type mismatch''. Looking at the Lean documentation, this behaviour occurs because the arguments inside curly brackets are not to be specified when calling the definition somewhere else in the code.
		
		Regarding error messages, I also found it particularly tricky to interpret them, as they appear (to my unexperienced eye) not strongly related to the true source of the error. In fact, Dr. Kevin Buzzard told us once that the Lean community knows about this issue and is actively seeking to make error messages more understandable.
		
		Formalising the results in \Cref{sec:informal} has also drawn my attention to details that I had overlooked before. For instance, a subtle point I realised is the interval where $f(x)$ is bounded on. When I was writing the informal proof I thought it would be enough to define what ``bounded on $(a,b)$'' is (see \cref{list:5}). However, when I was writing the proof of the main Lemma in Lean, I realised I needed to prove that $f$ is bounded on some neighbourhood $(x_0-\delta, x_0+\delta)$ of $x_0$, which is not necessarily $(a,b)$. Thus, I had to introduce a general definition of boundedness on an interval $I$ and then introduce the idea of $f$ being bounded on some interval $J\subseteq I$ (see \cref{list:7}). Obviously, in the Lean code we set $I:=(a,b)$.
		
		I also found it hard to find the right result to conclude the proof of the main Lemma in \Cref{sec:formal}. I had the hypothesis \lstinline|hltK1 : $\abs{f \langle x, hxDom \rangle} < K_1$| and \lstinline|$K_1$_le_$K_2$ : $K_1\leqslant K_2$| and I wanted to conclude that \lstinline|f $\abs{f \langle x, hxDom \rangle} \leqslant K_2$|. Since the \lstinline|exact?| tactic was not giving any useful results, I looked for the right lemma in the ``LeanSearch'' tool \cite{LeanSearch}, as one TA suggested in the labs. The correct tactic to use here is \lstinline|exact (le_of_lt hltK1).trans $K_1$_le_$K_2$|, where \lstinline|le_of_lt| gives the hypothesis \lstinline|$\abs{f \langle x, hxDom \rangle} \leqslant K_1$| and \lstinline|.trans| applies transitivity on non-strict inequalities.
		
		On a positive note, I found surprisingly easy writing the mathematical definitions. I also found useful the built-in search engine for lemmas in Lean (i.e., the \lstinline|exact?| and \lstinline|try?| instructions). Moreover, I found quite surprising the fact that the \lstinline|#lint| command picks up on missing docstrings for definitions or theorems. I added these to resolve the errors \lstinline|#lint| was showing at first.
		
		\subsection{Reflecting on the chosen lemma}
		
		The result on boundedness of functions around a point was not the one I initially intended to prove. At first I tried to prove the following:\\
		
		\begin{lemma*}{Finite Jensen's Inequality}
			Let $n\geq 2$ and $\varphi:\R\rightarrow\R$ be a convex function. If $\alpha_1,\dots,\alpha_n\in[0,1]$ satisfy $\alpha_1+\cdots+\alpha_n=1$, then the inequality
			\begin{equation*}
				\varphi(\alpha_1x_1+\cdots+\alpha_nx_n)\leq\alpha_1\varphi(x_1)+\cdots+\alpha_n\varphi(x_n)
			\end{equation*}
			holds for all $x_1,\dots,x_n\in\R$.
		\end{lemma*}
		
		This Lemma can be easily proved by induction on the number of points $n$. I broke down the proof in small lemmas that were easier by themselves. Once I had the small pieces formalised, however, I was unable to put together the induction step. I read about the \lstinline|induction'| tactic but was not able to use it satisfactorily to produce the final step.
		
	\section{Conclusions}\label{sec:conclu}

	In all, this has been an instructive introduction to Lean, where I have learnt to define concepts and to apply them to simple examples and easy lemmas. I have also proved the equivalence of two definitions of boundedness around a point, and further formally proved that if $f$ has a limit in $x_0$, then $f$ is bounded on some neighbourhood of $x_0$.
	
	For the upcoming projects I aim to prove some lemma or theorem that requires induction, since I did not succeed in my first attempt to prove Jensen's inequality. I am also interested in exploring other areas of mathematics such as Galois Theory or Complex Analysis, with an expectation to learn more tactics and discover the Mathlib library in greater depth.
	
	{\small
		\bibliographystyle{plainurl}
		\bibliography{references}
	}
		 
\end{document}